\chapter{Gauge-Invariant Einstein--Scalar Perturbations}

Scalar perturbations are reduced to a single gauge-invariant degree of freedom.  The continuum Mukhanov--Sasaki action is then discretized with the same projector Laplacian used elsewhere.

\allowdisplaybreaks

\section{Einstein--scalar input and background}
Let
\begin{equation}
M_{\mathrm{Pl}}^2:=(8\pi G)^{-1}
\end{equation}
and
\begin{equation}
S[g,\phi]
=\int\sqrt{-g}
\left[
\frac{M_{\mathrm{Pl}}^2}{2}R
-\frac12g^{\mu\nu}\partial_\mu\phi\partial_\nu\phi
-V(\phi)
\right]\dd^4x.
\end{equation}
For
\begin{equation}
\dd s^2=-\dd t^2+a(t)^2\dd\bm x^2,
\qquad
\phi=\bar\phi(t),
\end{equation}
the background equations are
\begin{align}
3M_{\mathrm{Pl}}^2H^2
&=\frac12\dot{\bar\phi}^{\,2}+V(\bar\phi),\\
-2M_{\mathrm{Pl}}^2\dot H
&=\dot{\bar\phi}^{\,2},\\
\ddot{\bar\phi}+3H\dot{\bar\phi}+V_{,\phi}
&=0,\\
\epsilon
&:=-\frac{\dot H}{H^2}
=\frac{\dot{\bar\phi}^{\,2}}
{2M_{\mathrm{Pl}}^2H^2}.
\end{align}

\section{Scalar gauge invariant}
Use the scalar convention
\begin{equation}
g_{ij}
=a^2\left[(1-2\Psi)\delta_{ij}
+2\partial_i\partial_jE\right].
\end{equation}
Under \(t\mapsto t+T\),
\begin{align}
\Psi&\mapsto\Psi+HT,\\
\delta\phi&\mapsto\delta\phi-\dot{\bar\phi}T.
\end{align}
Define
\begin{equation}
\boxed{
\zeta
:=-\Psi-\frac{H}{\dot{\bar\phi}}\delta\phi.
}
\end{equation}
Indeed, substitution of the two transformations gives
\begin{align}
\zeta
&\mapsto
-(\Psi+HT)
-\frac{H}{\dot{\bar\phi}}
\left(\delta\phi-\dot{\bar\phi}T\right)\\
&=-\Psi-HT
-\frac{H}{\dot{\bar\phi}}\delta\phi+HT\\
&=\zeta.
\end{align}
The time-shift contributions therefore cancel before any gauge choice is made.
In comoving gauge,
\begin{equation}
\delta\phi=0,
\qquad
E=0,
\qquad
h_{ij}=a^2\ee^{2\zeta}\delta_{ij}.
\end{equation}

\section{Lapse and shift elimination}
Write
\begin{equation}
N=1+\alpha,
\qquad
N_i=\partial_i\beta.
\end{equation}
The linear Hamiltonian and momentum constraints give
\begin{align}
\alpha&=\frac{\dot\zeta}{H},\\
\beta&=-\frac{\zeta}{H}+\chi,\\
\nabla^2\chi&=a^2\epsilon\dot\zeta.
\end{align}
Thus \(\alpha\) and \(\beta\) are nondynamical and are eliminated before the reduced quadratic action is formed.

\section{Reduced scalar action and equation}
After substitution of the constraints,
\begin{equation}
\boxed{
S_\zeta^{(2)}
=M_{\mathrm{Pl}}^2
\int a^3\epsilon
\left[
\dot\zeta^2-\frac{(\nabla\zeta)^2}{a^2}
\right]\dd t\dd^3x.
}
\end{equation}
For a compactly supported variation \(\delta\zeta\),
\begin{align}
\delta S_\zeta^{(2)}
&=2M_{\mathrm{Pl}}^2
\int
\left[
 a^3\epsilon\dot\zeta\,\delta\dot\zeta
-a\epsilon\partial_i\zeta\,\partial_i\delta\zeta
\right]\dd t\dd^3x\\
&=-2M_{\mathrm{Pl}}^2
\int
\left[
\frac{\dd}{\dd t}
\left(a^3\epsilon\dot\zeta\right)
-a\epsilon\nabla^2\zeta
\right]
\delta\zeta\,\dd t\dd^3x,
\end{align}
where the second line follows by integrating once in time and once in space.  Hence the Euler--Lagrange equation is
\begin{equation}
\boxed{
\frac{\dd}{\dd t}
\left(a^3\epsilon\dot\zeta\right)
-a\epsilon\nabla^2\zeta
=0.
}
\end{equation}
For a Fourier mode, \(\nabla^2\zeta_{\bm k}=-k^2\zeta_{\bm k}\), so
\begin{align}
0
&=\frac{\dd}{\dd t}
\left(a^3\epsilon\dot\zeta_{\bm k}\right)
+a\epsilon k^2\zeta_{\bm k}\\
&=a^3\epsilon\ddot\zeta_{\bm k}
+\left(3Ha^3\epsilon+a^3\dot\epsilon\right)
\dot\zeta_{\bm k}
+a\epsilon k^2\zeta_{\bm k}.
\end{align}
Division by \(a^3\epsilon\) gives
\begin{equation}
\ddot\zeta_{\bm k}
+\left(3H+\frac{\dot\epsilon}{\epsilon}\right)
\dot\zeta_{\bm k}
+\frac{k^2}{a^2}\zeta_{\bm k}=0.
\end{equation}
At \(k=0\), the equation can be integrated twice:
\begin{align}
\frac{\dd}{\dd t}
\left(a^3\epsilon\dot\zeta\right)&=0,\\
a^3\epsilon\dot\zeta&=C_2,\\
\dot\zeta&=\frac{C_2}{a^3\epsilon},\\
\zeta(t)&=C_1+C_2\int^t\frac{\dd s}{a(s)^3\epsilon(s)}.
\end{align}
Hence the adiabatic constant mode is conserved when the second mode decays.

\section{Mukhanov--Sasaki form}
Let
\begin{align}
\dd t&=a\dd\eta,\\
z(\eta)&:=a\sqrt{2\epsilon}\,M_{\mathrm{Pl}},\\
v&:=z\zeta.
\end{align}
Changing from \(t\) to \(\eta\) first gives
\begin{equation}
S_\zeta^{(2)}
=\frac12\int z^2
\left[(\zeta')^2-(\nabla\zeta)^2\right]
\dd\eta\dd^3x.
\end{equation}
Since \(\zeta=v/z\) and \(z\) is spatially homogeneous,
\begin{align}
z^2(\zeta')^2
&=\left(v'-\frac{z'}{z}v\right)^2\\
&=(v')^2
-\frac{z'}{z}(v^2)'
+\left(\frac{z'}{z}\right)^2v^2,\\
z^2(\nabla\zeta)^2&=(\nabla v)^2.
\end{align}
Integrating the mixed term by parts yields
\begin{align}
-\int\frac{z'}{z}(v^2)'\dd\eta
&=\int
\left(\frac{z'}{z}\right)'v^2\dd\eta
+\text{boundary term},\\
\left(\frac{z'}{z}\right)'
+\left(\frac{z'}{z}\right)^2
&=\frac{z''}{z}.
\end{align}
Therefore, up to the displayed boundary term,
\begin{equation}
\boxed{
S_v^{(2)}
=\frac12\int
\left[
(v')^2-(\nabla v)^2+\frac{z''}{z}v^2
\right]\dd\eta\dd^3x.
}
\end{equation}
Therefore
\begin{align}
v''-\nabla^2v-\frac{z''}{z}v&=0,\\
v_{\bm k}''+
\left(k^2-\frac{z''}{z}\right)v_{\bm k}&=0.
\end{align}

\section{Projector-lattice scalar action}
For cubic spacing \(h\),
\begin{equation}
S_{v,h}^{(2)}
=\frac12\int h^3\sum_{\bm j}
\left[
(v_{\bm j}')^2
-\sum_{r=1}^3(\delta_r^+v_{\bm j})^2
+\frac{z''}{z}v_{\bm j}^2
\right]\dd\eta.
\end{equation}
The discrete summation-by-parts identity
\begin{equation}
h^3\sum_{\bm j,r}
(\delta_r^+v_{\bm j})(\delta_r^+\delta v_{\bm j})
=-h^3\sum_{\bm j}
(\Delta_hv_{\bm j})\delta v_{\bm j}
\end{equation}
converts the gradient variation into the lattice Laplacian.  The lattice equation is therefore
\begin{equation}
v_{\bm j}''-\Delta_hv_{\bm j}-\frac{z''}{z}v_{\bm j}=0.
\end{equation}
For a lattice Fourier mode,
\begin{align}
\delta_r^+\ee^{\ii\bm k\cdot\bm x_{\bm j}}
&=\frac{\ee^{\ii k_rh}-1}{h}
\ee^{\ii\bm k\cdot\bm x_{\bm j}},\\
\left|\frac{\ee^{\ii k_rh}-1}{h}\right|^2
&=\frac{(\ee^{\ii k_rh}-1)(\ee^{-\ii k_rh}-1)}{h^2}\\
&=\frac{2-2\cos(k_rh)}{h^2}
=\frac4{h^2}\sin^2\left(\frac{k_rh}{2}\right).
\end{align}
Thus
\begin{align}
\widehat k_h^2
&:=\frac4{h^2}\sum_{r=1}^3
\sin^2\left(\frac{k_rh}{2}\right),\\
v_{\bm k,h}''+
\left(\widehat k_h^2-\frac{z''}{z}\right)v_{\bm k,h}&=0,\\
\widehat k_h^2
&=k^2-\frac{h^2}{12}\sum_{r=1}^3k_r^4
+\frac{h^4}{360}\sum_{r=1}^3k_r^6
+\order(h^6).
\end{align}
For \(z''/z=2/\eta^2\),
\begin{equation}
v_\kappa(\eta)
=\frac{1}{\sqrt{2\kappa}}
\left(1-\frac{\ii}{\kappa\eta}\right)
\ee^{-\ii\kappa\eta}
\end{equation}
solves
\begin{equation}
v_\kappa''+
\left(\kappa^2-\frac2{\eta^2}\right)v_\kappa=0.
\end{equation}
With \(\kappa=\widehat k_h\),
\begin{align}
p_{v,h}&=2.001500,\\
p_{|v|^2,h}&=2.002803,\\
p_{v,h}^{\mathrm{Richardson}}&=4.001170,\\
\max_h
\left|v_hv_h^{*\prime}-v_h^*v_h'-\ii\right|
&=3.219647\times10^{-15}.
\end{align}

\section{Open extensions}
\begin{nextsteps}
\item Extend the quadratic derivation from one canonical scalar to general background and transfer functions.
\item Add vector and tensor sources within the same gauge-invariant projector calculus.
\item Generalize to multiple scalar fields and derive the coupled adiabatic and entropy modes.
\item Derive the cubic action and the leading non-Gaussian correlation functions.
\item Obtain the vacuum normalization from the projector dynamics rather than prescribing it as an external input.
\end{nextsteps}
